\section{Related Work}

The computational demands of the JPEG compression standard \cite{jpeg_standard} have driven extensive research into hardware-accelerated Discrete Cosine Transform (DCT) architectures. While algorithmic optimizations like the canonical Loeffler fast DCT \cite{loeffler1989} minimize total arithmetic operations for software execution, mapping them directly to custom hardware often incurs irregular data flow and complex routing overhead.

Consequently, practical VLSI implementations typically favor row-column decomposition \cite{vlsi_dct_survey}. This approach computes the 2D DCT by sequentially applying a 1D DCT to the rows of an $8 \times 8$ block, passing the intermediate results through a transposition buffer, and applying a second 1D DCT to the columns. This decoupled structure drastically simplifies the datapath and control logic compared to true 2D matrix multipliers.

To maximize throughput within this row-column paradigm, hardware designs rely on two primary structural strategies: deep pipelining to increase the maximum operating frequency ($F_{max}$), and spatial parallelism (duplicating Multiply-Accumulate units) to process multiple pixels simultaneously. Rather than presenting a single static architecture, this work builds upon these foundational concepts to provide a highly parameterized SystemVerilog framework. By selectively enabling pipelining and 8-way parallelism, this project focuses heavily on systematically evaluating the empirical area, power, and performance tradeoffs of these techniques on a modern FPGA target.
